\section{Creation Awaiting Its Freedom}

The redemption of the body belongs to a promise wider than the human person considered alone. Paul describes creation groaning as it awaits freedom from corruption, and he joins its expectation to our own waiting for bodily redemption (Rom 8:18--25). The world God made is involved in the destiny of his children. The beauty of created things is real, their suffering and decay are real, and the final word about them belongs to the Creator. We cannot determine creation's consummation by extrapolating from the processes that now wear it away.

The new Jerusalem descends from God in Revelation 21. God dwells with his people, wipes away their tears, and makes all things new. Its center is his presence. The city has no need of a temple distinct from the Lord and the Lamb, for divine communion fills its life. The nations appear within this renewed world, bringing their glory and honor. Their common life is described through a dwelling received from God and through the goods the nations bring into it.

The same chapter excludes everything defiled. Renewal brings judgment upon evil as well as fulfillment of good. We cannot simply identify the future city with the continuation of whatever civilization happens to admire. Human works contain generosity and exploitation, beauty and vanity, truthful service and lies. Our work can serve the good promised by God, but this mixture needs his judgment and purification. God knows what in our labor belongs to charity and what must pass away. We serve best when we surrender both our work and our claim to control its ultimate form.

\emph{Gaudium et spes} 39 states that the manner and time of creation's transformation are unknown to us. It distinguishes earthly progress from the growth of Christ's kingdom, while teaching that hope for the new earth should strengthen our care for this one. The good fruits of human dignity, fraternity, and freedom, fostered in the Spirit, will be found purified and transfigured. We can therefore undertake responsible work for human dignity while entrusting its lasting fruit to God's transformation. Our labor is real and valuable; the final transfiguration is God's act.

That distinction sustains patience when good work appears fragile. A person can work for justice and still die before a wrong is repaired. The good already done for someone who was harmed retains its value, even while much remains unfinished. We are responsible for the fidelity of our service, and its final fruit belongs to God. The promise that charity endures lets us undertake imperfect work without either demanding immediate perfection or giving ourselves permission to do nothing.

Paul describes creation's expectation as groaning. We receive its goods amid real affliction. Bodies suffer, people injure one another, and places can be damaged by carelessness or greed. These injuries remain to be resisted and relieved while we wait for the redemption Paul promises. The final newness cannot be produced by multiplying our present arrangements; nevertheless, the love with which we tend a person or a place can already belong to the kingdom that is coming.

Waiting for creation's renewal gives us both work to do and something to entrust to God. Gratitude for its beauty need not become possessiveness, and sorrow over its wounds need not become despair. The Father who gives the world through his Son has not forgotten it. Our earthly home is passing in its present form, but its Maker is faithful. We can love what he has made while awaiting what he alone will make of it. When our part in a task ends, its unfinished good can still be entrusted to him.

